% Cristal de chlorure de sodium : empilement alterné d'ions sodium et chlorure (L'essentiel et exercice 1, chapitre 7)
\begin{tikzpicture}[x=1cm,y=1cm]
\foreach \k/\op in {2/0.35,1/0.65,0/1}{
  \begin{scope}[shift={({0.42*\k},{0.32*\k})},opacity=\op]
    \foreach \i in {0,1,2}{\draw[otMuted!60,line width=0.4pt] (0,{0.85*\i}) -- (1.7,{0.85*\i}) ({0.85*\i},0) -- ({0.85*\i},1.7);}
    \foreach \i in {0,1,2}{\foreach \j in {0,1,2}{
      \pgfmathtruncatemacro\parite{mod(\i+\j+\k,2)}
      \ifnum\parite=0
        \shade[ball color=otGreen!80] ({0.85*\i},{0.85*\j}) circle (0.24);
      \else
        \shade[ball color=otNavy!70] ({0.85*\i},{0.85*\j}) circle (0.13);
      \fi}}
  \end{scope}}
\shade[ball color=otGreen!80] (4.0,1.25) circle (0.24); \node[pcLegende,anchor=west] at (4.35,1.25) {ion chlorure $\ce{Cl-}$};
\shade[ball color=otNavy!70] (4.0,0.55) circle (0.13); \node[pcLegende,anchor=west] at (4.35,0.55) {ion sodium $\ce{Na+}$};
\end{tikzpicture}
\par\smallskip\legende{Chaque ion est entouré de six ions de charge opposée ; les interactions électrostatiques attractives assurent la cohésion du cristal.}
